\documentclass[11pt]{amsart}
\usepackage[localtheorems]{raeez-math-template}

\newtheorem{theorem}{Theorem}[section]
\newtheorem{proposition}[theorem]{Proposition}
\newtheorem{lemma}[theorem]{Lemma}
\newtheorem{corollary}[theorem]{Corollary}
\newtheorem{hypothesis}[theorem]{Hypothesis}
\theoremstyle{definition}
\newtheorem{definition}[theorem]{Definition}
\newtheorem{example}[theorem]{Example}
\theoremstyle{remark}
\newtheorem{remark}[theorem]{Remark}

\providecommand{\At}{\mathrm{At}}
\providecommand{\BCOV}{\mathrm{BCOV}}
\providecommand{\CoHA}{\operatorname{CoHA}}
\providecommand{\End}{\mathrm{End}}
\providecommand{\Etwo}{\mathsf{E}_2}
\providecommand{\Ethree}{\mathsf{E}_3}
\providecommand{\Eone}{\mathsf{E}_1}
\providecommand{\HH}{\mathrm{HH}}
\providecommand{\HKR}{\mathrm{HKR}}
\providecommand{\MC}{\mathrm{MC}}
\providecommand{\Perf}{\mathrm{Perf}}
\providecommand{\PV}{\mathrm{PV}}
\providecommand{\QME}{\mathrm{QME}}
\providecommand{\Sym}{\mathrm{Sym}}
\providecommand{\Td}{\mathrm{Td}}
\providecommand{\ch}{\mathrm{ch}}
\providecommand{\cA}{\mathcal A}
\providecommand{\cC}{\mathcal C}
\providecommand{\cL}{\mathcal L}
\providecommand{\cO}{\mathcal O}
\providecommand{\cW}{\mathcal W}
\providecommand{\dbar}{\overline{\partial}}
\providecommand{\fg}{\mathfrak g}
\providecommand{\fgl}{\mathfrak{gl}}
\providecommand{\hCS}{\mathrm{hCS}}
\providecommand{\id}{\mathrm{id}}
\providecommand{\kanom}{\kappa_{\mathrm{anom}}}
\providecommand{\kBKM}{\kappa_{\mathrm{BKM}}}
\providecommand{\kcat}{\kappa_{\mathrm{cat}}}
\providecommand{\kch}{\kappa_{\mathrm{ch}}}
\providecommand{\kfiber}{\kappa_{\mathrm{fiber}}}
\providecommand{\C}{\mathbb C}
\providecommand{\Z}{\mathbb Z}

\title{HKR, Atiyah Classes, and BCOV Curving on Calabi-Yau Threefolds}
\author{Raeez Lorgat}
\date{2026-04-30}

\begin{document}
\maketitle

\begin{abstract}
For a smooth complex variety, Hochschild chains and cochains admit the
Hochschild-Kostant-Rosenberg comparison with forms and polyvectors
after the Hochschild complex is interpreted on the completed diagonal.
On affine space the comparison is the elementary symbol map. On a global
smooth variety the product, module, and Mukai-pairing comparisons carry
the Kontsevich-Caldararu Duflo correction, equivalently the square-root
Todd class generated by the Atiyah class. This correction belongs to
Hochschild calculus and formality. Kapranov's Atiyah-class brackets,
BCOV closed-string curving, Positselski curvature, and the open
holomorphic Chern-Simons quantum master equation live in separate
deformation problems. Matrix hCS fields occupy the open BV layer.
Hochschild polyvectors and Calabi-Yau categorical data occupy two
separate algebraic layers. CoHA inputs belong to the Hall enumerative
layer. The compact Calabi-Yau threefold theory requires the full
Costello quantization package: orientation and framing, regulator,
gauge fixing, completions, anomaly trivialization, and factorization or
TCFT gluing.
\end{abstract}

\section{Hochschild and Polyvector Calculus}
\label{sec:hkr-calculus}

Let $X$ be a smooth complex variety and let $\Perf(X)$ be its dg
category of perfect complexes. Hochschild cochains and chains are
written
\[
  \HH^\bullet(X)=\HH^\bullet(\Perf(X)), \qquad
  \HH_\bullet(X)=\HH_\bullet(\Perf(X)).
\]
The Dolbeault polyvector and Hodge complexes are
\[
  \PV^\bullet(X)=
  \bigoplus_{p,q}\Omega^{0,q}(X,\wedge^pT_X),
  \qquad
  \Omega^\bullet(X)=
  \bigoplus_{p,q}\Omega^{0,q}(X,\Omega^p_X).
\]

\begin{hypothesis}[HKR and completed diagonal]
\label{hyp:hkr-diagonal}
The HKR statements below use one of the following smooth
settings: a smooth separated finite-type scheme over $\C$, a smooth
complex manifold with perfect coherent input, or a formal affine chart
with continuous Hochschild complexes. Hochschild chains and cochains are
computed for $\Perf(X)$ by kernels on $X\times X$ completed along the
diagonal $\Delta\colon X\hookrightarrow X\times X$:
\[
  \widehat{X\times X}_{\Delta}
  =
  \varprojlim_n
  \operatorname{Spec}_{X\times X}
  \bigl(\cO_{X\times X}/\mathcal I_\Delta^{n+1}\bigr).
\]
All tensor products and endomorphisms on this formal neighbourhood are
$\mathcal I_\Delta$-adically complete. Properness enters exactly for
finite-dimensional traces, Mukai pairings, and global partition
functions.
\end{hypothesis}

For smooth $X$ there are HKR quasi-isomorphisms
\[
  I_{\HKR}^{\mathrm{co}}\colon \HH^\bullet(X)\simeq
  \bigoplus_{p,q}H^q(X,\wedge^pT_X),
  \qquad
  I_{\HKR}^{\mathrm{ho}}\colon \HH_\bullet(X)\simeq
  \bigoplus_{p,q}H^q(X,\Omega^p_X).
\]
They identify the associated graded objects for the Hodge filtration.
The affine formula is the antisymmetrisation map
\[
  \xi_1\wedge\cdots\wedge\xi_p
  \longmapsto
  \frac{1}{p!}\sum_{\sigma\in S_p}\mathrm{sgn}(\sigma)\,
  \xi_{\sigma(1)}\otimes\cdots\otimes\xi_{\sigma(p)}
\]
from polyvector fields to Hochschild cochains.

\begin{definition}[Atiyah class]
\label{def:atiyah-class}
For a holomorphic vector bundle $E$ on $X$, the Atiyah class
\[
  \At(E)\in H^1(X,\Omega^1_X\otimes\End(E))
\]
is the extension class of the first-jet sequence
\[
  0\to\Omega^1_X\otimes E\to J^1(E)\to E\to0.
\]
For $E=T_X$, this class is the first obstruction to a holomorphic
connection on $T_X$ and is the source of Kapranov's brackets on
$T_X[-1]$.
\end{definition}

\begin{theorem}[Corrected HKR comparison]
\label{thm:corrected-hkr}
Let $X$ satisfy Hypothesis~\ref{hyp:hkr-diagonal}. The uncorrected HKR
maps are quasi-isomorphisms of complexes and identify the associated
graded Hochschild calculus with polyvectors and forms. Compatibility
with the Hochschild module action on chains, the Mukai pairing, and the
global formality descent is given by the Duflo-Todd correction. On
homology the corrected map is
\[
  I_{\HKR,\Td}^{\mathrm{ho}}(\alpha)
  =
  I_{\HKR}^{\mathrm{ho}}(\alpha)\wedge \Td(T_X)^{1/2}.
\]
On cochains the same formal-coordinate descent acts by the Duflo
element of the Lie algebra of formal vector fields. In Chern-Weil
language its characteristic class is generated by the Atiyah class
through the exponential series
\[
  \operatorname{Duf}_X
  =
  \exp\!\left(
    \sum_{m\ge1}
    \frac{B_{2m}}{4m(2m)!}\,
    \operatorname{tr}\bigl(\At(T_X)^{2m}\bigr)
  \right),
\]
whose image on Hochschild homology is $\Td(T_X)^{1/2}$.
The corrected comparison is the Kontsevich-Caldararu, Calaque-Van den
Bergh, and Caldararu-Willerton form of HKR.
\end{theorem}

\begin{proof}
HKR gives the local symbol quasi-isomorphism. Kontsevich formality
identifies polyvector fields and polydifferential operators as
$L_\infty$ algebras on formal affine charts. Global descent over the
bundle of formal coordinates introduces the Gelfand-Fuks characteristic
class represented by $\At(T_X)$. The Duflo element of this descent is
the square-root Todd class. Caldararu's Mukai-pairing theorem and the
Calaque-Van den Bergh cochain comparison identify this same class as
the correction required for product and module compatibility.
\end{proof}

\begin{corollary}[Flat affine three-space]
\label{cor:c3-hkr}
For $X=\C^3$ the tangent bundle is globally trivial, hence
\[
  \At(T_{\C^3})=0,\qquad \Td(T_{\C^3})=1.
\]
The corrected HKR comparison equals the elementary HKR symbol map:
\[
  \HH^\bullet(\C^3)\simeq
  \C[z_1,z_2,z_3]\otimes
  \wedge(\partial_{z_1},\partial_{z_2},\partial_{z_3})
  =\PV^\bullet(\C^3).
\]
This is the flat calculation. It supplies the local model for hCS
Feynman rules and for the positive CoHA of $\C^3$.
\end{corollary}

\section{Formality, Atiyah Brackets, and Curvature}
\label{sec:formality-curvature}

The symbols HKR, formality, Atiyah class, BCOV curving, and
Positselski curvature refer to different structures.

\begin{definition}[Formal diagonal]
\label{def:formal-diagonal}
The formal diagonal of a smooth $X$ is the completed neighbourhood
$\widehat{X\times X}_{\Delta}$ of
Hypothesis~\ref{hyp:hkr-diagonal}. Its structure sheaf carries the
completed jet algebra
\[
  \cJ_X=\Delta^{-1}\cO_{\widehat{X\times X}_{\Delta}},
  \qquad
  \operatorname{gr}\cJ_X\simeq \widehat{\Sym}_{\cO_X}(\Omega^1_X).
\]
Splitting the first infinitesimal layer of $\cJ_X$ is the same first
order problem as splitting the jet sequence of $T_X$; the obstruction
class is $\At(T_X)$. The formal diagonal is the geometric home of HKR
descent and Kapranov's shifted Lie construction.
\end{definition}

\begin{hypothesis}[Kapranov shifted Lie data]
\label{hyp:kapranov-shifted-lie}
Assume $X$ is a smooth complex manifold. Choose a $C^\infty$
connection $\nabla$ on $T_X$ whose curvature representative has
Dolbeault class $\At(T_X)$. The completed Dolbeault complex
\[
  \Omega^{0,\bullet}(X,T_X[-1])^{\wedge}
\]
is taken at the zero section. Kapranov's brackets are the multilinear
maps $\ell_n$ obtained from $R_\nabla$ and its iterated covariant
derivatives:
\[
  \ell_2 \longleftrightarrow [R_\nabla]=\At(T_X),
  \qquad
  \ell_n \longleftrightarrow \nabla^{n-2}R_\nabla
  \quad (n\ge3),
\]
with the usual symmetrisation and contraction conventions. Different
choices of $\nabla$ give $L_\infty$-isomorphic completed objects.
\end{hypothesis}

\begin{definition}[Kontsevich-Caldararu formality]
\label{def:kc-formality}
Kontsevich formality is an $L_\infty$ quasi-isomorphism from
polyvector fields to polydifferential operators. Caldararu's global
form applies this comparison to smooth varieties by inserting the
Atiyah-Duflo descent class. The output is a corrected Hochschild
calculus comparison.
\end{definition}

\begin{definition}[Kapranov brackets]
\label{def:kapranov-brackets}
Under Hypothesis~\ref{hyp:kapranov-shifted-lie}, Kapranov's
construction gives $T_X[-1]$ the structure of a Lie algebra object in
the derived category. In a Dolbeault model this is an $L_\infty$
algebra whose binary bracket is induced by $\At(T_X)$. The higher
brackets are iterated covariant derivatives and contractions of the
curvature class representing $\At(T_X)$. This structure controls the
formal neighbourhood of the diagonal. hCS quantization and chiral
transfer add BV and factorization data beyond this formal moduli
problem.
\end{definition}

\begin{definition}[BCOV closed theory]
\label{def:bcov-theory}
For a Calabi-Yau threefold $(X,\Omega_X)$, BCOV theory is the closed
Kodaira-Spencer theory on polyvectors with Calabi-Yau trace. Its
tree-level cubic on complex-structure directions is the Yukawa form
\[
  Y_X(v_1,v_2,v_3)
  =
  \int_X
  \Omega_X\wedge
  \bigl((v_1\wedge v_2\wedge v_3)\,\lrcorner\,\Omega_X\bigr),
  \qquad v_i\in H^1(X,T_X).
\]
The Costello-Li one-loop BCOV datum is a determinant-line anomaly or
closed BV curving term. In the standard compact threefold
normalisation its scalar coefficient contains $\chi_{\mathrm{top}}(X)/24$.
\end{definition}

\begin{definition}[Positselski curvature]
\label{def:positselski-curvature}
A curved dg algebra in Positselski's sense is a triple $(A,d,h)$ with
$d^2=[h,-]$ and $d(h)=0$. The element $h$ is algebraic curvature for
coderived or contraderived categories. It is a curvature element inside
an algebraic differential. It is separate from the BCOV closed BV
curving and from the Atiyah class.
\end{definition}

\begin{proposition}[The six roles]
\label{prop:six-roles}
The following assignments are type-correct.
\[
\begin{array}{c|c|c}
\text{object} & \text{home} & \text{role}\\
\hline
I_{\HKR} & \HH^\bullet(X)\to\PV^\bullet(X) & \text{quasi-isomorphism}\\
\Td(T_X)^{1/2} & H^{\mathrm{even}}(X) & \text{HKR product and pairing correction}\\
\At(T_X) & H^1(X,\Omega^1_X\otimes\End T_X) & \text{Kapranov binary bracket source}\\
Y_X & \Sym^3H^1(X,T_X)^\vee & \text{tree-level BCOV cubic}\\
\alpha_{\BCOV} & \text{closed BV anomaly line} & \text{one-loop BCOV datum}\\
h & A^2 & \text{Positselski curved-dg curvature}
\end{array}
\]
The Atiyah class can source several of these constructions through
descent or trace, yet the resulting objects live in different
cohomological groups and answer different deformation problems.
\end{proposition}

\begin{proof}
The HKR map is a comparison between Hochschild and polyvector
complexes. The Todd class is a characteristic-class correction to this
comparison. Kapranov's bracket is a Lie-object structure on $T_X[-1]$.
The Yukawa cubic is a scalar trilinear form after contraction with the
Calabi-Yau volume form. The BCOV one-loop datum belongs to the closed
BV quantization. Positselski curvature is an element controlling the
square of an algebraic differential. These domains are functorially
connected, but none has the same source, target, or degree as another.
\end{proof}

\begin{remark}[Complex deformation unobstructedness]
\label{rem:btt}
For a compact Kahler Calabi-Yau manifold, the Kodaira-Spencer dg Lie
algebra governing complex deformations is homotopy abelian by the
Bogomolov-Tian-Todorov theorem. This statement concerns the deformation
functor of complex structures. Kapranov's Atiyah brackets and the BCOV
Yukawa cubic remain meaningful cohomological operations. The Kuranishi
obstruction map may vanish while the tree-level Yukawa coupling remains
nonzero.
\end{remark}

\section{Four Local Layers}
\label{sec:layers}

Let $X$ be a Calabi-Yau threefold with holomorphic volume form
$\Omega_X$, and let $V$ be a finite-dimensional vector space.

\begin{definition}[Open hCS field theory]
\label{def:hcs-fields}
The matrix holomorphic Chern-Simons fields on an open set $U\subset X$
are
\[
  \cA\in\Omega^{0,\bullet}(U,\fgl(V))[1].
\]
The classical action is
\[
  S_{\hCS}(\cA)
  =
  \int_U\Omega_X\wedge
  \operatorname{Tr}_V\left(
    \frac12\cA\,\dbar\cA+
    \frac16\cA[\cA,\cA]\right).
\]
\end{definition}

\begin{proposition}[Four layers]
\label{prop:four-layers}
For $X=\C^3$ with matrix coefficient $V$, the following layers are
distinct.
\[
\begin{array}{c|c|c}
\text{layer} & \text{datum} & \text{operation}\\
\hline
\text{Calabi-Yau category} & \Perf(\C^3),\ \operatorname{tr}_{\Omega} &
  \text{cyclic trace and Serre pairing}\\
\text{Hochschild polyvectors} & \HH^\bullet(\C^3)\simeq\PV^\bullet(\C^3) &
  \text{Gerstenhaber and formality calculus}\\
\text{matrix hCS} & \Omega^{0,\bullet}(\C^3,\fgl(V))[1] &
  \text{open BV action}\\
\text{CoHA} & H^\bullet_{\mathrm{eq}}(\mathcal M_{\C^3},\phi_W) &
  \text{Hall multiplication on critical sheaf moduli}
\end{array}
\]
The matrix coefficient $\End(V)$ supplies the open gauge sector. The
CoHA is built separately from moduli of objects and vanishing cycles.
\end{proposition}

\begin{proof}
The first row is Morita-invariant categorical data. The second row is
the Hochschild calculus of that category. The third row is a local BV
field theory whose fields are Dolbeault forms with values in
$\fgl(V)$. The fourth row is an enumerative Hall algebra. On affine
space the HKR symbol map and Costello's local hCS observables compare
the first three rows after the flat Duflo correction collapses. The
Hall multiplication uses correspondences between moduli stacks and is
an additional construction.
\end{proof}

\begin{theorem}[The affine CoHA normalisation]
\label{thm:c3-coha}
For the toric Calabi-Yau threefold $\C^3$,
\[
  \CoHA(\C^3)=Y^+(\widehat{\mathfrak{gl}}_1),
\]
the positive half of the affine Yangian of $\mathfrak{gl}_1$. A full
$\cW_{1+\infty}$ algebra is obtained only after a double, centre,
completion, or evaluation construction has been specified.
\end{theorem}

\begin{proof}
The cohomological Hall algebra of the one-loop three-arrow quiver with
potential $xyz-xzy$ gives the positive Hall half. Schiffmann-Vasserot
and Maulik-Okounkov identify this positive half with the positive part
of the affine Yangian of $\mathfrak{gl}_1$. The full
$\cW_{1+\infty}$ requires both positive and negative modes together
with the central and completion data.
\end{proof}

\section{Compact hCS Quantization}
\label{sec:compact-hcs}

\begin{hypothesis}[Compact hCS package]
\label{hyp:compact-hcs}
A compact hCS quantization statement on a Calabi-Yau threefold $X$
requires the following data.
\begin{enumerate}[label=\textup{(\roman*)}]
\item a holomorphic volume form $\Omega_X$ and an orientation of the
determinant line defining the BV pairing;
\item a quadratic dg Lie algebra $(\fg,\langle-,-\rangle)$ and, when
colour traces are used, a representation $\rho:\fg\to\End(W)$;
\item an elliptic gauge fixing $Q^{\mathrm{GF}}$, heat-kernel
regularisation, propagators, and renormalisation counterterms;
\item completed local functionals and observables on which the BV
Laplacian, bracket, and RG flow are defined;
\item a solution of the scale-dependent quantum master equation;
\item vanishing or explicit trivialisation of the quartic anomaly
class in the Costello local obstruction complex;
\item factorization descent over $X$ and TCFT sewing data for global
traces or partition functions.
\end{enumerate}
\end{hypothesis}

\begin{proposition}[Conditional compact observable algebra]
\label{prop:compact-observables}
Under Hypothesis~\ref{hyp:compact-hcs}, the completed quantum hCS
observables form a quantum factorization algebra on $X$. If the
quartic obstruction class survives in the local BV obstruction complex
for the chosen fields, regulator, and counterterms, the compact quantum
observable algebra has no value at that order.
\end{proposition}

\begin{proof}
Costello's renormalisation constructs the effective interaction
$I[L]$ scale by scale. The quantum factorization algebra exists exactly
when the RG flow is compatible with the BV operator and $I[L]$
satisfies the QME at every scale. The obstruction to extending a
partial solution lies in the local BV cohomology. A surviving class
prevents the next-scale effective interaction from satisfying the QME.
\end{proof}

\begin{definition}[Quartic invariant]
\label{def:quartic-invariant}
For a representation $\rho:\fg\to\End(W)$, define
\[
  I_4^\rho(x_1,x_2,x_3,x_4)
  =
  \frac1{4!}\sum_{\sigma\in S_4}
  \operatorname{Tr}_W\!\left(
    \rho(x_{\sigma(1)})\rho(x_{\sigma(2)})
    \rho(x_{\sigma(3)})\rho(x_{\sigma(4)})\right).
\]
Then $I_4^\rho\in\Sym^4(\fg^\vee)^\fg$.
\end{definition}

\begin{proposition}[QME anomaly]
\label{prop:qme-anomaly}
In the flat Costello local model in complex dimension three, after a
choice of gauge fixing and regulator, the nonabelian one-loop open hCS
obstruction to the QME is represented by
\[
  \Theta^\rho_U(\alpha_0,\alpha_1,\alpha_2,\alpha_3)
  =
  \int_U
  \left[
  I_4^\rho\bigl(
    \alpha_0,\partial\alpha_1,\partial\alpha_2,\partial\alpha_3
  \bigr)\right]_{(3,3)},
  \qquad \alpha_i\in\Omega_c^{0,\bullet}(U,\fg).
\]
Its class is
\[
  \kanom(X,\fg,\rho)
  =
  [\,\hbar\,\Theta^\rho\,]\in
  H^1_{\mathrm{loc}}
  \bigl(\Omega^{0,\bullet}(X,\fg)[1],\cO_{\mathrm{loc}}\bigr).
\]
\end{proposition}

\begin{proof}
The contributing one-loop graph is the four-external-leg wheel. Its
colour factor is the symmetrised trace of four representation matrices,
which is $I_4^\rho$. The holomorphic form factor contains three
$\partial$ derivatives, producing a top holomorphic density in complex
dimension three. The Costello obstruction complex places the resulting
local functional in degree one, precisely the degree of the first QME
obstruction.
\end{proof}

\begin{definition}[Quadratic Casimir]
\label{def:quadratic-casimir}
Let $\{T_a\}$ be a basis of $\fg$ and let $\{T^a\}$ be the dual basis
for the invariant pairing. The representation-theoretic quadratic
Casimir is
\[
  C_2^\rho=\sum_a\rho(T_a)\rho(T^a).
\]
For a simple Lie algebra in the adjoint normalisation,
$C_2^{\mathrm{ad}}=2h^\vee$.
\end{definition}

\begin{proposition}[Kinetic counterterm]
\label{prop:kinetic-counterterm}
With the flat heat-kernel normalisation, the two-point graph produces
a degree-zero kinetic counterterm
\[
  S_{\mathrm{ct}}^{(2)}
  =
  -\hbar\,\frac{C_2^\rho}{(4\pi)^3}
  \log(L/\varepsilon)
  \int_X\Omega_X\wedge\operatorname{Tr}_\rho(\cA\,\dbar\cA),
\]
up to the chosen pairing and regulator convention. The quartic class of
Proposition~\ref{prop:qme-anomaly} is the QME obstruction. The
quadratic Casimir counterterm changes the kinetic representative.
\end{proposition}

\begin{proof}
Closing two legs in the cubic hCS vertex yields the colour contraction
$\sum_a\rho(T_a)\rho(T^a)=C_2^\rho$. The two remaining external fields
have the tensor type of the kinetic functional. The short-distance
six-real-dimensional heat-kernel integral supplies the logarithmic
factor. Its cohomological degree is zero, while the QME obstruction has
degree one.
\end{proof}

\section{The Native Operadic Level of \texorpdfstring{$\Phi_3$}{Phi3}}
\label{sec:phi3-level}

\begin{hypothesis}[Chiral and factorization transfer]
\label{hyp:chiral-transfer}
Fix a Calabi-Yau category $\cC$ of dimension $d$, a smooth CY target
$X$, a holomorphic factorization algebra $\Phi^{\mathrm{FA}}_d(\cC)$
constructed from the Hochschild calculus with the chosen formality
datum, and a witnessed admissible specialisation datum
$(\Sigma_{d-1},C)$. The specialisation
\[
  \Sp^{\mathrm{ch}}_{\Sigma_{d-1},C}
  \bigl(\Phi^{\mathrm{FA}}_d(\cC)\bigr)
\]
uses completed observables, proper-support or compact-support
conditions along $\Sigma_{d-1}$, descent for the factorization
structure, and sewing on the reference curve $C$. For compact hCS
inputs it also uses Hypothesis~\ref{hyp:compact-hcs}.
\end{hypothesis}

\begin{proposition}[Stage and output]
\label{prop:stage-output}
Under Hypothesis~\ref{hyp:chiral-transfer}, for a Calabi-Yau category
of dimension $d\ge3$, the native chiral output of the functor $\Phi_d$
is $\Eone$. Higher operadic structures appear before specialisation or
after passage to centres, doubles, and module categories.
\end{proposition}

\begin{proof}
The Hochschild and factorization-algebra input of a $d$-Calabi-Yau
category carries the expected $E_d$ or shifted Gerstenhaber structure.
The chiral specialisation step to a curve lowers the native output in
dimension $d\ge3$ to an $\Eone$ chiral algebra. An $\Etwo$ structure on
the same data requires a separate centre, double, or module-layer
construction such as $Z(\operatorname{Rep}(A))$ or a Drinfeld double.
\end{proof}

\section{\texorpdfstring{$K3\times E$}{K3 x E}}
\label{sec:k3e}

Let $X=K3\times E$, with $E$ an elliptic curve.

\begin{theorem}[Five invariant readings]
\label{thm:k3e-readings}
The standard invariant readings on $K3\times E$ are
\[
\begin{aligned}
  \kcat(K3\times E)
  &=\chi(\cO_{K3})\chi(\cO_E)=2\cdot0=0,\\
  \kch(A_{K3\times E})
  &=\sum_q(-1)^qh^{0,q}(K3\times E)=0,\\
  \kch^{\mathrm{Heis}}(K3\times E)
  &=2+1=3,\\
  \kBKM(\mathfrak g_{\Delta_5})
  &=c_1(0)/2=10/2=5,\\
  \kfiber(K3)&=\operatorname{rk}H^\bullet(K3,\Z)=24.
\end{aligned}
\]
They arise from different constructions: categorical Euler
characteristic, compact Hodge supertrace, Heisenberg-Cartan
specialisation, Borcherds product weight, and Mukai-lattice fibre rank.
Cross-construction sums are undefined.
\end{theorem}

\begin{proof}
Kunneth for coherent cohomology gives
\[
  \chi(\cO_{K3\times E})=
  \chi(\cO_{K3})\chi(\cO_E)=2\cdot0=0.
\]
The compact chiral Hodge supertrace is the same alternating
$(0,q)$-column:
\[
  \sum_q(-1)^qh^{0,q}(K3\times E)
  =
  (h^{0,0}(K3)-h^{0,1}(K3)+h^{0,2}(K3))
  (h^{0,0}(E)-h^{0,1}(E))
  =
  (1+1)(1-1)=0.
\]
The Heisenberg-Cartan branch is rank-additive over the product-stable
specialisation and gives the K3 contribution $2$ plus the elliptic
free-boson contribution $1$. The Borcherds value is the weight of the
Gritsenko product $\Delta_5$, equivalently $c_1(0)/2=10/2=5$. The
fibre value $24$ is the rank of $H^\bullet(K3,\Z)$ with the Mukai
lattice.
\end{proof}

\begin{remark}[Atiyah and BCOV on the product]
\label{rem:k3e-atiyah-bcov}
The elliptic tangent bundle is holomorphically trivial, so
$\At(T_E)=0$. The product decomposition gives the Atiyah class of
$T_{K3\times E}$ from the K3 summand. The scalar one-loop BCOV Euler
coefficient vanishes because
\[
  \chi_{\mathrm{top}}(K3\times E)=\chi_{\mathrm{top}}(K3)\chi_{\mathrm{top}}(E)
  =24\cdot0=0.
\]
The Todd factor in HKR descent still contains the K3 term
$1+c_2(K3)/24$. The vanishing of the Euler coefficient supplies no
global strict HKR trivialisation by itself.
\end{remark}

\section{References}
\label{sec:references}

\begin{itemize}[leftmargin=1.2em,itemsep=2pt]
\item M. F. Atiyah, \emph{Complex analytic connections in fibre
bundles}, Trans. AMS 85 (1957), 181-207.
\item R. Hochschild, B. Kostant, and A. Rosenberg,
\emph{Differential forms on regular affine algebras}, Trans. AMS 102
(1962), 383-408.
\item M. Kontsevich, \emph{Deformation quantization of Poisson
manifolds}, Lett. Math. Phys. 66 (2003), 157-216.
\item M. Kontsevich, \emph{Deformation quantization of algebraic
varieties}, EuroConf. Mosh Flato 2000, volume 1, 2000.
\item A. Caldararu, \emph{The Mukai pairing II: the Hochschild-Kostant-
Rosenberg isomorphism}, Adv. Math. 194 (2005), 34-66.
\item D. Calaque and M. Van den Bergh, \emph{Hochschild cohomology and
Atiyah classes}, Adv. Math. 224 (2010), 1839-1889.
\item A. Caldararu and S. Willerton, \emph{The Mukai pairing I: a
categorical approach}, New York J. Math. 16 (2010), 61-98.
\item M. Kapranov, \emph{Rozansky-Witten invariants via Atiyah
classes}, Compositio Math. 115 (1999), 71-113.
\item P. Deligne, P. Griffiths, J. Morgan, and D. Sullivan,
\emph{Real homotopy theory of Kahler manifolds}, Invent. Math. 29
(1975), 245-274.
\item G. Tian, \emph{Smoothness of the universal deformation space of
compact Calabi-Yau manifolds and its Petersson-Weil metric},
Mathematical Aspects of String Theory, 1987.
\item A. Todorov, \emph{The Weil-Petersson geometry of the moduli space
of SU(n) Calabi-Yau manifolds}, Comm. Math. Phys. 126 (1989), 325-346.
\item M. Bershadsky, S. Cecotti, H. Ooguri, and C. Vafa,
\emph{Kodaira-Spencer theory of gravity and exact results for quantum
string amplitudes}, Comm. Math. Phys. 165 (1994), 311-427.
\item P. Candelas, X. de la Ossa, P. Green, and L. Parkes,
\emph{A pair of Calabi-Yau manifolds as an exactly soluble
superconformal theory}, Nucl. Phys. B 359 (1991), 21-74.
\item K. Costello, \emph{Topological conformal field theories and
Calabi-Yau categories}, Adv. Math. 210 (2007), 165-214.
\item K. Costello, \emph{Renormalization and Effective Field Theory},
AMS Mathematical Surveys and Monographs 170, 2011.
\item K. Costello and S. Li, \emph{Quantum BCOV theory on Calabi-Yau
manifolds and the higher genus B-model}, 2012.
\item K. Costello and S. Li, \emph{Twisted supergravity and its
quantization}, 2016.
\item K. Costello and O. Gwilliam, \emph{Factorization Algebras in
Quantum Field Theory}, Volumes 1 and 2.
\item L. Positselski, \emph{Two kinds of derived categories, Koszul
duality, and comodule-contramodule correspondence}, Mem. AMS 212
(2011).
\item O. Schiffmann and E. Vasserot, \emph{The elliptic Hall algebra
and the K-theory of the Hilbert scheme of A2}, Duke Math. J. 162
(2013), 279-366.
\item A. Negut, \emph{The shuffle algebra revisited}, Int. Math. Res.
Notices 2014, 6242-6275.
\end{itemize}

\end{document}
